\section{Mid-Term Examination Suite}

\ExamHeader{Mid-Term Examination Suite}{Set C (Model Paper 3)}{90 Minutes}{30 Marks}{-0.25 Marks}

\noindent
\textbf{General Instructions:}
\begin{enumerate}[topsep=2pt,itemsep=1pt]
    \item This paper contains \textbf{30 Multiple-Choice Questions (MCQs)} carrying \textbf{1 mark each} (Total = 30 Marks).
    \item \textbf{10 Questions from Unit 1} (EM Theory), \textbf{10 Questions from Unit 2} (Lasers), and \textbf{10 Questions from Unit 3} (Fiber Optics).
    \item Each incorrect response incurs a deduction of \textbf{0.25 marks}.
\end{enumerate}
\vspace{3mm}

\subsection*{Section A: Unit 1 --- Electromagnetic Theory (Questions 1 to 10)}
\mcqitem{1}{The divergence of the electric field $\nabla \cdot \vec{E}$ in a medium of charge density $\rho$ is:}{$\rho/\epsilon$}{0}{$-\rho/\epsilon$}{$\mu_0 J$}
\mcqitem{2}{If $\vec{B} = \nabla \times \vec{A}$, then $\nabla \cdot \vec{B}$ is:}{0}{1}{$\vec{A}$}{$\rho$}
\mcqitem{3}{The scalar electric potential $V$ satisfies $\vec{E} = $}{$-\nabla V$}{$\nabla V$}{$\nabla \times V$}{$\nabla^2 V$}
\mcqitem{4}{In an electromagnetic wave in free space, $\vec{E}$ and $\vec{B}$ are:}{Perpendicular to each other and to direction of propagation}{Parallel}{At $45^\circ$}{Opposite}
\mcqitem{5}{The equation $\nabla^2 V = 0$ is known as:}{Laplace's equation}{Poisson's equation}{Ampere's law}{Faraday's law}
\mcqitem{6}{The continuity equation expresses conservation of:}{Electric charge}{Energy}{Momentum}{Mass}
\mcqitem{7}{The unit of electric permittivity $\epsilon_0$ is:}{$\text{F/m}$}{$\text{H/m}$}{$\text{N/C}$}{$\text{Tesla}$}
\mcqitem{8}{Displacement current flows through:}{Dielectric between capacitor plates during charging}{Copper wire}{Resistor}{Battery}
\mcqitem{9}{Stokes' theorem relates:}{Line integral and surface integral}{Volume and surface}{Scalar and vector}{None}
\mcqitem{10}{The Poynting vector magnitude represents:}{Electromagnetic power density ($W/m^2$)}{Force}{Electric charge}{Voltage}

\subsection*{Section B: Unit 2 --- Lasers and Applications (Questions 11 to 20)}
\mcqitem{11}{Stimulated emission produces photons that are:}{Coherent with incident photon}{Incoherent}{Random in frequency}{Opposite in phase}
\mcqitem{12}{Population inversion is achieved when:}{$N_2 > N_1$}{$N_1 > N_2$}{$N_1 = N_2$}{$N_2 = 0$}
\mcqitem{13}{Einstein's $A_{21}$ coefficient represents the probability of:}{Spontaneous emission}{Stimulated emission}{Absorption}{Pumping}
\mcqitem{14}{The wavelength emitted by He-Ne laser is:}{$632.8\text{ nm}$}{$1060\text{ nm}$}{$900\text{ nm}$}{$532\text{ nm}$}
\mcqitem{15}{The pumping source in a Ruby laser is:}{Helical Xenon flash lamp}{Electric discharge}{Chemical reaction}{Thermal}
\mcqitem{16}{Ruby laser is a:}{3-level laser}{4-level laser}{2-level laser}{Semiconductor laser}
\mcqitem{17}{In a He-Ne laser, the active medium that produces lasing is:}{Neon}{Helium}{Argon}{Krypton}
\mcqitem{18}{Semiconductor lasers require:}{Direct bandgap material (GaAs)}{Indirect bandgap (Si)}{Gas tube}{Flash lamp}
\mcqitem{19}{Holography was discovered by:}{Dennis Gabor}{Maiman}{Einstein}{Townes}
\mcqitem{20}{A hologram records:}{Both amplitude and phase}{Amplitude only}{Phase only}{Wavelength only}

\subsection*{Section C: Unit 3 --- Fiber Optics (Questions 21 to 30)}
\mcqitem{21}{Light propagates through optical fiber by:}{Total Internal Reflection}{Diffraction}{Refraction}{Interference}
\mcqitem{22}{In an optical fiber, the refractive indices satisfy:}{$n_1 > n_2$ (Core > Cladding)}{$n_1 < n_2$}{$n_1 = n_2$}{$n_2 = 1$}
\mcqitem{23}{The Numerical Aperture of an optical fiber is:}{$NA = \sqrt{n_1^2 - n_2^2}$}{$NA = n_1 - n_2$}{$NA = n_1 + n_2$}{$NA = n_1/n_2$}
\mcqitem{24}{The cutoff V-number for single-mode fiber is:}{$V \le 2.405$}{$V > 2.405$}{$V = 10$}{$V = 0$}
\mcqitem{25}{In a graded-index fiber, refractive index decreases:}{Radially from center to cladding}{Axially}{Linearly along length}{Randomly}
\mcqitem{26}{The primary cause of intrinsic attenuation in silica glass at 1550 nm is:}{Rayleigh scattering ($\propto 1/\lambda^4$)}{Absorption}{Bending}{Splice loss}
\mcqitem{27}{The lowest optical loss window in silica fibers is at:}{$1550\text{ nm}$ (~0.2 dB/km)}{$850\text{ nm}$}{$1310\text{ nm}$}{$632.8\text{ nm}$}
\mcqitem{28}{In endoscopy, image transmission is performed using a:}{Coherent fiber bundle}{Non-coherent bundle}{Single thick glass rod}{Mirror}
\mcqitem{29}{Attenuation $\alpha$ in dB/km for length $L$ is given by:}{$\alpha = \frac{10}{L}\log_{10}(P_{\text{in}}/P_{\text{out}})$}{$\alpha = \frac{P_{\text{in}}}{P_{\text{out}}}$}{$\alpha = \ln(P_{\text{in}}/P_{\text{out}})$}{$\alpha = 10 L$}
\mcqitem{30}{The number of modes in a graded-index fiber is approximately:}{$V^2/4$}{$V^2/2$}{$V$}{$2V$}

\newpage
\subsection*{Complete Answer Key \& Technical Rationales}

\begin{table}[htbp]
\centering
\caption{\textbf{Mid-Term Answer Key \& Explanations (Questions 1 to 30)}}
\vspace{2mm}
\begin{tabularx}{\textwidth}{l c X l c X}
\toprule
\textbf{Q\#} & \textbf{Ans} & \textbf{Technical Rationale} & \textbf{Q\#} & \textbf{Ans} & \textbf{Technical Rationale} \\
\midrule
1 & \textbf{(a)} & Gauss's law in differential form. & 16 & \textbf{(a)} & Maiman's 3-level solid-state ruby laser. \\
2 & \textbf{(a)} & Divergence of curl is identically zero. & 17 & \textbf{(a)} & Neon atoms emit laser light; Helium transfers collision energy. \\
3 & \textbf{(a)} & Electric field is the negative gradient of electrostatic potential. & 18 & \textbf{(a)} & Direct bandgap allows efficient radiative recombination. \\
4 & \textbf{(a)} & Electromagnetic waves are strictly transverse waves. & 19 & \textbf{(a)} & Invented by Dennis Gabor in 1948. \\
5 & \textbf{(a)} & Laplace's equation in electrostatic charge-free regions. & 20 & \textbf{(a)} & Complete wavefront reconstruction. \\
6 & \textbf{(a)} & Fundamental charge conservation. & 21 & \textbf{(a)} & Continuous TIR along core-cladding interface. \\
7 & \textbf{(a)} & Farads per meter (or $C^2/(N\cdot m^2)$). & 22 & \textbf{(a)} & Core must be optically denser. \\
8 & \textbf{(a)} & Time-varying electric flux generates displacement current in dielectrics. & 23 & \textbf{(a)} & Defines light-gathering power of the fiber. \\
9 & \textbf{(a)} & Line integral of vector equals surface integral of curl. & 24 & \textbf{(a)} & Condition for only fundamental $HE_{11}$ mode propagation. \\
10 & \textbf{(a)} & Directional energy flux density vector. & 25 & \textbf{(a)} & Parabolic profile equalizes modal transit times. \\
11 & \textbf{(a)} & Stimulated emission creates identical, in-phase photons. & 26 & \textbf{(a)} & Rayleigh scattering scales as $1/\lambda^4$. \\
12 & \textbf{(a)} & Excited state population exceeds lower state population. & 27 & \textbf{(a)} & Optimal wavelength for long-haul telecommunications. \\
13 & \textbf{(a)} & $A_{21}$ is spontaneous transition rate per second. & 28 & \textbf{(a)} & Coherent bundle preserves spatial position of pixels. \\
14 & \textbf{(a)} & Standard visible red laser line. & 29 & \textbf{(a)} & Standard decibel loss definition. \\
15 & \textbf{(a)} & Optical pumping via flash tube. & 30 & \textbf{(a)} & Graded index supports half the modes of step index ($V^2/4$). \\
\bottomrule
\end{tabularx}
\end{table}
